\documentclass[12pt]{article}
\def\activityid{F07-U-v3}\usepackage[letterpaper,margin=.65in,top=.60in,bottom=.65in]{geometry}
\usepackage[T1]{fontenc}\usepackage[default]{sourcesanspro}
\usepackage{microtype,tikz,array,tabularx,fancyhdr,amsmath}\usepackage[hidelinks]{hyperref}
\usetikzlibrary{arrows.meta,calc}
\setlength{\parindent}{0pt}\setlength{\parskip}{7pt}\setlength{\headheight}{0pt}\setlength{\footskip}{26pt}\setlength{\emergencystretch}{2em}
\pagestyle{fancy}\fancyhf{}\renewcommand{\headrulewidth}{0pt}\renewcommand{\footrulewidth}{.4pt}
\fancyfoot[L]{\scriptsize Bellingham Math Circle / Week 7 / \activityid}\fancyfoot[R]{\scriptsize \thepage}
\definecolor{ink}{gray}{.12}\definecolor{soft}{gray}{.95}\color{ink}
\newcommand{\PageTitle}[2]{{\fontsize{10}{12}\selectfont\bfseries WEEK 7\quad GAME LAB\hfill #1\par}\vspace{3pt}{\fontsize{26}{29}\selectfont\bfseries #2\par}\vspace{2pt}}
\newcommand{\Subhead}[1]{\par\vspace{3pt}{\large\bfseries #1}\par}
\newcommand{\RuleBox}[1]{\par\begingroup\setlength{\fboxsep}{8pt}\colorbox{soft}{\parbox{\dimexpr\linewidth-16pt\relax}{#1}}\endgroup\par}
\newcommand{\WriteBox}[2]{\par\textbf{#1}\par\vspace{2pt}\begin{tikzpicture}\draw[black!40,line width=.5pt] (0,0) rectangle (\dimexpr\linewidth-.5pt\relax,#2);\end{tikzpicture}\par}
\newcommand{\Code}[1]{\texttt{#1}}

\newcommand{\StudentHeader}[1]{{\fontsize{10}{12}\selectfont\bfseries Week 7 / Strategy lab / #1\par}\vspace{10pt}}
\newcommand{\Problem}[2]{\par\vspace{4pt}\textbf{Problem #1:} #2\par}
\newcommand{\RecordBox}[2]{\par #1\par\vspace{2pt}\begin{tikzpicture}\draw[black!40,line width=.5pt](0,0)rectangle(\dimexpr\linewidth-.5pt\relax,#2);\end{tikzpicture}\par}

\hypersetup{pdftitle={Week 7: Grades 4--5}}
\begin{document}
\StudentHeader{Grades 4--5}
\Problem{1}{Two players alternate taking 1, 3, or 4 counters. Taking 2 is not allowed. Never take more than remain. Taking the last counter wins. Play from 8, 10, and 14. Record your moves. Swap who starts for a second attempt at a pile that surprised you.}
\begin{center}\renewcommand{\arraystretch}{2.1}\begin{tabular}{c|c|p{2.2in}|p{1.4in}}Starting pile & First move & Next moves, in order & Winner: first or second\\\hline
8 &  &  & \\
10 &  &  & \\
14 &  &  & \\
 & & & \\\end{tabular}\end{center}
\Problem{2}{W means the next player can force a win; L means every legal move gives the opponent a winning strategy. Zero is L. Build small piles and work upward: mark W if a legal move reaches L, and L if every legal move reaches W. Record the next piles as well as the labels.}
\begin{center}\renewcommand{\arraystretch}{1.55}\begin{tabular}{c|p{3.5in}|c}Pile & Legal next piles & W/L\\\hline
0 & none & L\\
1 &  & \\
2 &  & \\
3 &  & \\
4 &  & \\
5 &  & \\
6 &  & \\
7 &  & \\\end{tabular}\end{center}

\newpage
\StudentHeader{Grades 4--5}
\Problem{3}{Continue using moves 1, 3, and 4. Fill the chart in order. In each cell write L, or write W and a winning move: W1 means W with winning move 1. Check all legal next piles before marking L. Compare the rows after filling at least two of them.}
\begin{center}\renewcommand{\arraystretch}{1.9}\begin{tabular}{p{.9in}|*{7}{>{\centering\arraybackslash}p{.55in}|}}Pile & 0 & 1 & 2 & 3 & 4 & 5 & 6\\\hline
Label/move &  &  &  &  &  &  & \\\hline
Pile & 7 & 8 & 9 & 10 & 11 & 12 & 13\\
Label/move &  &  &  &  &  &  & \\\hline
Pile & 14 & 15 & 16 & 17 & 18 & 19 & 20\\
Label/move &  &  &  &  &  &  & \\\hline
Pile & 21 & 22 & 23 & 24 & 25 & 26 & 27\\
Label/move &  &  &  &  &  &  & \\\hline\end{tabular}\end{center}
\Problem{4}{Compare these four piles. Record every legal next pile and its label. Circle the starting piles that have no move to L.}
\begin{center}\renewcommand{\arraystretch}{1.9}\begin{tabular}{c|p{1.65in}|p{1.65in}|p{1.65in}}Pile & Take 1: next pile/label & Take 3: next pile/label & Take 4: next pile/label\\\hline
14 &  &  & \\
16 &  &  & \\
15 &  &  & \\
17 &  &  & \\\end{tabular}\end{center}
\RecordBox{The repeating pattern I see:}{0.55in}

\newpage
\StudentHeader{Grades 4--5}
\Problem{5}{Arrange 14, 16, 15, and 17 counters in whole groups of 7 with a smaller leftover group. Record the leftovers. Match each pile to its label in Problem 4.}
\begin{center}\renewcommand{\arraystretch}{1.7}\begin{tabular}{c|c|c|c}Pile & Whole groups of 7 & Leftover counters & W/L\\\hline
14 &  &  & \\
16 &  &  & \\
15 &  &  & \\
17 &  &  & \\\end{tabular}\end{center}
\Problem{6}{Use your chart to choose a winning move for each leftover count marked W. Record the leftover count after the move. For counts marked L, list every leftover count an allowed move can reach; use piles 7 and 9 so that all three moves are possible.}
\begin{center}\renewcommand{\arraystretch}{1.9}\begin{tabular}{c|c|p{2.3in}|p{1.7in}}Leftover count & W/L & Move to L, or all moves from L & New leftover count(s)\\\hline
0 &  &  & \\
1 &  &  & \\
2 &  &  & \\
3 &  &  & \\
4 &  &  & \\
5 &  &  & \\
6 &  &  & \\\end{tabular}\end{center}
\Problem{7}{Build 24 counters. Choose a first move using your table, then play. After each of your turns record the leftover count when grouping the remaining pile in sevens.}
\RecordBox{First move and later leftover counts:}{0.6in}

\newpage
\StudentHeader{Grades 4--5}
\Problem{8}{Change the allowed moves to 1, 3, and 5. Play from 6 and from 7, then swap who starts. Build smaller piles to check a proposed strategy. Record which piles you want to hand the other player.}
\RecordBox{My games and the piles I want to leave:}{0.9in}
\Problem{9}{Return to moves 1, 3, and 4. Copy the labels for piles 0--3 into the first row below and 7--10 into the second. These are four-label windows. To find the next label, inspect the first, second, and fourth cells: mark the next W if any is L, and L if all three are W. Record the next label, then slide the window one cell right.}
\begin{center}\renewcommand{\arraystretch}{2.0}\begin{tabular}{c|p{1.4in}|c|p{1.6in}}Window starts at & Four labels & Next label & Window after sliding\\\hline
0 &  &  & \\
7 &  &  & \\
1 &  &  & \\
8 &  &  & \\\end{tabular}\end{center}
\Problem{10}{Find another matching pair of four-label windows in your chart. Copy both windows, extend them by two labels, and compare the results.}
\RecordBox{My matching windows and their next two labels:}{0.9in}

\end{document}
